\documentclass[a4paper,11pt]{article}
% Preamble: same conventions as notes/equilibrium-energy-v3/preamble.tex
% (single-document build, no standalone chapter wrappers).
\usepackage[T2A]{fontenc}
\usepackage[utf8]{inputenc}
\usepackage[russian]{babel}
\usepackage{amsmath}
\usepackage{amssymb}

\numberwithin{equation}{section}
\setlength{\emergencystretch}{2em}

\newcommand{\rf}{\mathrm{rf}}
\newcommand{\eq}{\mathrm{eq}}

\usepackage[hidelinks]{hyperref}


\title{Хроматичность кольца: от гамильтониана до краевого поля}
\author{}
\date{}

\begin{document}
\maketitle

\begin{abstract}
Заметка собирает в одном месте определение хроматичности и её вычисление. Хроматичность определена через одно-оборотное отображение на дисперсионной замкнутой орбите и через усреднение гамильтониана по бетатронной фазе; показано, что оба определения совпадают. Хроматичность разложена на слагаемые по их происхождению в гамильтониане: кинематический (квадрупольный) член, кривизна тела диполя, перекрёстный член с $D_1'$, секступоли и вклад жёсткого края диполя, который следует из $\operatorname{rot}\mathbf{B}=0$ и в модели <<только тело>> отсутствует. Разобран переход от конечного краевого поля к жёсткому краю. Описано, какие из этих слагаемых учитывают BELL, COSY INFINITY, MAD-X, PTC и OptiM, и как хроматичность извлекается при наличии ВЧ. Численных результатов в заметке нет.
\end{abstract}

\tableofcontents

\clearpage
\section{Язык задачи}
\label{sec:setup}

В этой главе введены переменные и гамильтониан, на которых построено всё дальнейшее. Независимая переменная --- длина $s$ вдоль референсной орбиты, поперечные импульсы нормированы на импульс синхронной частицы. Отклонение энергии задаётся тремя способами --- через импульс, полную и кинетическую энергию; приведены формулы перевода хроматичности между ними, поскольку разные программы пользуются разными переменными.

\subsection{Обозначения}

Индекс $s$ отмечает синхронную (референсную) частицу.

\begin{center}
\renewcommand{\arraystretch}{1.2}
\begin{tabular}{@{}lp{0.55\textwidth}@{}}
\hline
\multicolumn{2}{@{}l}{\emph{Синхронная частица}}\\
$p_s$, $E_s=\gamma_s mc^2$, $K_s$ & импульс, полная и кинетическая энергия\\
$\beta_s$, $\gamma_s$ & скорость в единицах $c$ и лоренц-фактор\\
$C_s$ & длина орбиты\\
\hline
\multicolumn{2}{@{}l}{\emph{Оптика}}\\
$1/\rho$ & кривизна референсной орбиты (в диполе), $0$ вне диполей\\
$K_1$, $K_2$ & градиент квадруполя и секступоля, нормированные на $B\rho=p_s/q$\\
$K_x=K_1+1/\rho^2$, $K_y=-K_1$ & жёсткости линейных уравнений\\
$\alpha$, $\beta$, $\gamma$ & функции Твисса (по $x$ или $y$)\\
$D_1$, $D_1'$ & дисперсия на $\delta$ и её производная\\
$\nu_{x,y}$, $\mu=2\pi\nu$ & числа бетатронных колебаний и набег фазы за оборот\\
$J$, $\psi$ & инвариант Куранта--Снайдера и бетатронная фаза\\
$\langle\cdot\rangle_\psi$ & среднее по бетатронной фазе\\
\hline
\multicolumn{2}{@{}l}{\emph{Отклонение энергии}}\\
$\delta$ & $=\delta_p$, относительное отклонение импульса \eqref{eq:dp}\\
$\delta_E$, $\delta_K$ & то же для полной и кинетической энергии \eqref{eq:dE-dK}\\
$\xi=\mathrm{d}\nu/\mathrm{d}\delta$ & хроматичность\\
\hline
\multicolumn{2}{@{}l}{\emph{Резонатор}}\\
$h$, $\omega_\rf$, $V$ & кратность, частота, амплитуда напряжения\\
\hline
\end{tabular}
\end{center}

Буква $h$ закреплена за кратностью ВЧ; кривизна всюду пишется как $1/\rho$.

\subsection{Гамильтониан в сопровождающей системе}

Сопровождающая система координат связана с референсной орбитой: $x$ --- отклонение в плоскости поворота (наружу от центра кривизны), $y$ --- вертикальное. Канонические импульсы $p_x$, $p_y$ нормированы на $p_s$; вне векторного потенциала $p_x=P_x/p_s$, где $P_x$ --- механический импульс. Гамильтониан есть взятая с минусом $s$-компонента канонического импульса, делённая на $p_s$. Элемент дуги на расстоянии $x$ от орбиты длиннее $\mathrm{d}s$ в $1+x/\rho$ раз, поэтому кинетическая часть несёт этот множитель. В диполе с однородным полем $B_0$, $qB_0/p_s=1/\rho$, векторный потенциал в сопровождающих координатах сводится к полиному. Вместе с полями квадруполя и секступоля
\begin{equation}
\label{eq:H}
\mathcal{H}
=-\Bigl(1+\frac{x}{\rho}\Bigr)\sqrt{(1+\delta)^2-p_x^2-p_y^2}
+\frac{x}{\rho}+\frac{x^2}{2\rho^2}
+\frac{K_1}{2}\bigl(x^2-y^2\bigr)
+\frac{K_2}{6}\bigl(x^3-3xy^2\bigr).
\end{equation}
Поля в \eqref{eq:H} фиксированы: от $\delta$ зависит только кинетическая часть. Это главное свойство канонической формулировки --- вся хроматичность происходит из корня и множителя $1+x/\rho$.

В линейном порядке уравнения Гамильтона дают
\begin{equation}
\label{eq:hill}
x''+\frac{K_x}{1+\delta}\,x=\frac{\delta}{\rho(1+\delta)},
\qquad
y''+\frac{K_y}{1+\delta}\,y=0,
\end{equation}
т.\,е. при импульсе $p_s(1+\delta)$ жёсткости делятся на $1+\delta$.

Если ВЧ нет, $\delta$ --- сохраняющийся параметр, и движение четырёхмерно. С ВЧ добавляется пара (продольная координата, $\delta$); в COSY это $(l,\delta_K)$ с $l=-(t-t_0)v_0\gamma_s/(1+\gamma_s)$~\cite{ValetovPhD}. В статическом поле вместо $\delta$ можно взять пару $(t,-E)$, как в~\cite{EqNote}. До главы~\ref{sec:rf} ВЧ нет.

\subsection{Переменные энергии и перевод хроматичности}

Определим
\begin{equation}
\label{eq:dp}
\delta=\delta_p=\frac{p-p_s}{p_s},
\end{equation}
\begin{equation}
\label{eq:dE-dK}
\delta_E=\frac{E-E_s}{E_s},
\qquad
\delta_K=\frac{K-K_s}{K_s},
\qquad
p_t=\frac{E-E_s}{p_sc}.
\end{equation}
Последняя переменная --- \texttt{PT} в MAD-X~\cite{MADX}. Из $\mathrm{d}E=v\,\mathrm{d}p$ и $\mathrm{d}K=\mathrm{d}E$ в первом порядке
\begin{equation}
\label{eq:conv}
\delta_E=\beta_s^2\,\delta,
\qquad
p_t=\beta_s\,\delta,
\qquad
\delta_K=\Bigl(1+\frac{1}{\gamma_s}\Bigr)\delta,
\end{equation}
где использовано $E_s\beta_s^2/K_s=(\gamma_s^2-1)/(\gamma_s(\gamma_s-1))=(\gamma_s+1)/\gamma_s$. Хроматичность --- производная, поэтому
\begin{equation}
\label{eq:xi-conv}
\xi\equiv\frac{\mathrm{d}\nu}{\mathrm{d}\delta}
=\Bigl(1+\frac{1}{\gamma_s}\Bigr)\frac{\mathrm{d}\nu}{\mathrm{d}\delta_K}
=\beta_s\frac{\mathrm{d}\nu}{\mathrm{d}p_t}
=\beta_s^2\frac{\mathrm{d}\nu}{\mathrm{d}\delta_E}.
\end{equation}
Перевод точен для первой производной в точке $\delta=0$. Для нелинейной хроматичности \eqref{eq:conv} нужно брать с поправками следующих порядков (например, $\delta_K$ содержит член $\propto\delta^2$), и коэффициенты при $\delta^j$, $j\ge2$, переводятся не простым умножением. Ниже везде $\xi=\mathrm{d}\nu/\mathrm{d}\delta$ при $\delta=0$ и бесконечно малой амплитуде, натуральная, если не оговорено иное.

\section{Что такое хроматичность}
\label{sec:def}

Хроматичность --- производная числа бетатронных колебаний по $\delta$. Числа колебаний частицы с $\delta\neq0$ считаются около её собственной замкнутой орбиты, которая сдвинута на дисперсию. В этой главе дано два определения: через одно-оборотное отображение (след матрицы в неподвижной точке) и через гамильтониан, усреднённый по бетатронной фазе. Показано, что они совпадают. В конце введена нелинейная хроматичность.

\subsection{Дисперсионная замкнутая орбита}

При $\delta\neq0$ замкнутая орбита частицы
\begin{equation}
\label{eq:co}
x_{\mathrm{co}}(s;\delta)=D_1(s)\,\delta+D_2(s)\,\delta^2+O(\delta^3),
\qquad
p_{x,\mathrm{co}}=D_1'\,\delta+O(\delta^2),
\end{equation}
где $D_1$ --- периодическое решение уравнения \eqref{eq:hill} в первом порядке:
\begin{equation}
\label{eq:D1}
D_1''+K_x D_1=\frac{1}{\rho}.
\end{equation}
Бетатронное движение отсчитывается от этой орбиты: $x=x_{\mathrm{co}}+X$, $p_x=p_{x,\mathrm{co}}+P_x$, $y=Y$, $p_y=P_y$ (вертикальной дисперсии в плоском кольце нет).

\subsection{Определение через одно-оборотное отображение}

Пусть $\mathcal{M}_\delta$ --- одно-оборотное отображение поперечных переменных $z=(x,p_x,y,p_y)$ при фиксированном $\delta$. Замкнутая орбита есть неподвижная точка
\begin{equation}
\label{eq:fixed}
\mathcal{M}_\delta(\mathring z)=\mathring z,
\qquad
\mathring z(\delta)=\bigl(D_1,D_1',0,0\bigr)\delta+O(\delta^2).
\end{equation}
Линеаризация около неё --- матрица $\tilde M(\delta)=\partial\mathcal{M}_\delta/\partial z\,|_{\mathring z}$. Для плоского несвязанного кольца она блочно-диагональна, и
\begin{equation}
\label{eq:nu-trace}
\cos\mu_{x,y}(\delta)=\tfrac12\operatorname{tr}\tilde M_{x,y}(\delta),
\qquad
\nu=\frac{\mu}{2\pi},
\end{equation}
причём ветвь $\arccos$ выбирается по знаку элемента $\tilde M_{12}$. Отсюда
\begin{equation}
\label{eq:xi-trace}
\xi_{x,y}=-\frac{1}{4\pi\sin\mu_{x,y}}\,
\frac{\mathrm{d}}{\mathrm{d}\delta}\operatorname{tr}\tilde M_{x,y}\Big|_{\delta=0}.
\end{equation}
Производная следа складывается из двух частей: прямой зависимости отображения от $\delta$ и сдвига точки линеаризации:
\begin{equation}
\label{eq:dM}
\frac{\mathrm{d}\tilde M}{\mathrm{d}\delta}
=\frac{\partial^2\mathcal{M}}{\partial z\,\partial\delta}
+\frac{\partial^2\mathcal{M}}{\partial z^2}\cdot\frac{\partial\mathring z}{\partial\delta}.
\end{equation}
Второе слагаемое --- квадратичные коэффициенты отображения, умноженные на дисперсию. Если его опустить (линеаризовать около референсной орбиты), секступоли и геометрическая нелинейность диполя из хроматичности выпадут. Явная запись \eqref{eq:xi-trace} через коэффициенты отображения дана в приложении~\ref{app:map}. Именно так хроматичность получают из карт Тейлора~\cite{ValetovPhD,ValetovChrom2026}.

\subsection{Определение через усреднённый гамильтониан}

Подставим $x=D_1\delta+X$, $p_x=D_1'\delta+P_x$ в \eqref{eq:H} и разложим до первого порядка по $\delta$. Часть, квадратичная по бетатронным переменным, имеет вид
\begin{equation}
\label{eq:H2}
\mathcal{H}_2=\mathcal{H}_2^{(0)}+\delta\,\mathcal{H}_2^{(1)}+O(\delta^2),
\qquad
\mathcal{H}_2^{(1)}=\tfrac12\,Z^{\mathsf T}G(s)\,Z,
\quad Z=(X,P_x)\ \text{или}\ (Y,P_y).
\end{equation}
Члены, линейные по $Z$, с $\delta$ уничтожаются выбором $D_1$ \eqref{eq:D1}. Хроматичность есть сдвиг набега фазы под действием $\delta\,\mathcal{H}_2^{(1)}$:
\begin{equation}
\label{eq:xi-ham}
\xi=\frac{1}{2\pi}\oint\frac{\partial\langle\mathcal{H}_2^{(1)}\rangle_\psi}{\partial J}\,\mathrm{d}s
\equiv\frac{1}{4\pi}\oint I\,\mathrm{d}s,
\end{equation}
со средними по фазе
\begin{equation}
\label{eq:avg}
\langle X^2\rangle_\psi=J\beta,
\qquad
\langle XP\rangle_\psi=-J\alpha,
\qquad
\langle P^2\rangle_\psi=J\gamma.
\end{equation}
Множитель в определении $I$ выбран так, что добавка $\Delta K\,X^2/2$ даёт привычное $\Delta\nu=(1/4\pi)\oint\beta\,\Delta K\,\mathrm{d}s$. Формула \eqref{eq:xi-ham} --- то же определение, что в~\cite{EqNote}, \S\,2.3.

\subsection{Эквивалентность}

Покажем, что \eqref{eq:xi-ham} совпадает с \eqref{eq:xi-trace}. Обозначим $A(s)=\begin{pmatrix}\beta&-\alpha\\-\alpha&\gamma\end{pmatrix}$ и $\mathcal{J}=\begin{pmatrix}0&1\\-1&0\end{pmatrix}$. Тогда $\langle ZZ^{\mathsf T}\rangle_\psi=JA$ по \eqref{eq:avg}, а одно-оборотная матрица от точки $s$ есть
\begin{equation}
\label{eq:Ms}
M_s=I\cos\mu+A\mathcal{J}\sin\mu.
\end{equation}
Возмущение $\delta\,\mathcal{H}_2^{(1)}$ на отрезке $\mathrm{d}s$ меняет движение как $Z'=\mathcal{J}\,\partial\mathcal{H}/\partial Z$, т.\,е. добавляет к матрице отрезка $\delta\,\mathcal{J}G\,\mathrm{d}s$. В первом порядке
\begin{equation}
\label{eq:dtr}
\begin{aligned}
\mathrm{d}\operatorname{tr}\tilde M
&=\delta\oint\operatorname{tr}\bigl(M_s\,\mathcal{J}G\bigr)\,\mathrm{d}s
=\delta\oint\Bigl[\cos\mu\operatorname{tr}(\mathcal{J}G)+\sin\mu\operatorname{tr}(A\mathcal{J}\mathcal{J}G)\Bigr]\mathrm{d}s\\
&=-\delta\sin\mu\oint\operatorname{tr}(AG)\,\mathrm{d}s,
\end{aligned}
\end{equation}
так как $\operatorname{tr}(\mathcal{J}G)=0$ для симметричной $G$ и $\mathcal{J}^2=-I$. С другой стороны $\mathrm{d}\operatorname{tr}\tilde M=-2\sin\mu\,\mathrm{d}\mu$ и $\langle\mathcal{H}_2^{(1)}\rangle_\psi=\tfrac12\operatorname{tr}(G\langle ZZ^{\mathsf T}\rangle)=\tfrac{J}{2}\operatorname{tr}(AG)$. Значит,
\begin{equation}
\label{eq:equiv}
\mathrm{d}\mu=\frac{\delta}{2}\oint\operatorname{tr}(AG)\,\mathrm{d}s
=\delta\oint\frac{\partial\langle\mathcal{H}_2^{(1)}\rangle_\psi}{\partial J}\,\mathrm{d}s,
\end{equation}
что и есть \eqref{eq:xi-ham}. Неподвижная точка \eqref{eq:fixed} в гамильтоновом языке --- это подстановка дисперсионной орбиты, а производная следа --- среднее по фазе. Оба определения используют один и тот же гамильтониан; если в нём есть краевые члены (глава~\ref{sec:edge}), они войдут в обе формулы одинаково, через тонкую <<линзу>> $G\propto\delta_D(s-s_0)$.

\subsection{Нелинейная хроматичность}

Зависимость $\nu(\delta)$ при конечном $\delta$ записывается рядом
\begin{equation}
\label{eq:nu-series}
\nu(\delta)=\sum_{j\ge0}\frac{\nu^{(j)}}{j!}\,\delta^j,
\qquad
\xi=\nu^{(1)},
\qquad
\xi^{(j)}=\nu^{(j+1)},
\end{equation}
как в~\cite{ValetovChrom2026}. Определение \eqref{eq:nu-trace} остаётся в силе при любом $\delta$; гамильтонов путь для $\nu^{(2)}$ требует дисперсии второго порядка $D_2$, членов $\mathcal{H}_4$ и второго порядка теории возмущений. На практике высшие $\nu^{(j)}$ берут из карт высокого порядка (глава~\ref{sec:codes}).

\section{Из чего складывается хроматичность}
\label{sec:terms}

Разложение гамильтониана \eqref{eq:H} около дисперсионной орбиты даёт $\mathcal{H}_2^{(1)}$ в виде суммы нескольких слагаемых разного происхождения. Каждое из них после усреднения \eqref{eq:avg} даёт свой вклад в $I_{x,y}$ и в $\xi_{x,y}$. В этой главе перечислены вклады тела элементов: кинематический (он же квадрупольный), кривизна тела диполя, перекрёстный член с $D_1'$ и секступоли. Вклад края диполя отложен до главы~\ref{sec:edge}.

\subsection{Разложение гамильтониана}

До третьего порядка по $(x,p_x,p_y,\delta)$ корень в \eqref{eq:H} равен $(1+\delta)-(p_x^2+p_y^2)/(2(1+\delta))$, и
\begin{equation}
\label{eq:H3}
\mathcal{H}\simeq
-\frac{x\,\delta}{\rho}
+\Bigl(1+\frac{x}{\rho}\Bigr)\frac{p_x^2+p_y^2}{2(1+\delta)}
+\frac{K_x}{2}x^2+\frac{K_y}{2}y^2
+\frac{K_2}{6}\bigl(x^3-3xy^2\bigr)
\end{equation}
(без постоянной). Первый член --- источник дисперсии, $K_xx^2/2$ и $K_yy^2/2$ от $\delta$ не зависят. Подстановка $x=D_1\delta+X$, $p_x=D_1'\delta+P_x$ и сбор членов $\delta\times(\text{квадратичные по }X,P_x,Y,P_y)$ дают
\begin{equation}
\label{eq:H21}
\mathcal{H}_2^{(1)}
=\underbrace{-\frac{P_x^2+P_y^2}{2}}_{\mathrm{A}}
+\underbrace{\frac{D_1}{\rho}\,\frac{P_x^2+P_y^2}{2}}_{\mathrm{C}}
+\underbrace{\frac{D_1'}{\rho}\,X P_x}_{\mathrm{E}}
+\underbrace{\frac{K_2D_1}{2}\bigl(X^2-Y^2\bigr)}_{\mathrm{S}}.
\end{equation}
Член A --- от множителя $1/(1+\delta)$, C --- от множителя $1+x/\rho$ при сдвиге $x$ на $D_1\delta$, E --- от того же множителя при сдвиге $p_x$ на $D_1'\delta$ (произведение $X/\rho$ на $2D_1'\delta P_x$), S --- от кубического члена секступоля.

\subsection{Вклады по членам}

По \eqref{eq:xi-ham}, \eqref{eq:avg}
\begin{equation}
\label{eq:I-all}
I_x=-\gamma_x+\frac{D_1}{\rho}\gamma_x-\frac{2D_1'}{\rho}\alpha_x+K_2D_1\beta_x,
\qquad
I_y=-\gamma_y+\frac{D_1}{\rho}\gamma_y-K_2D_1\beta_y,
\end{equation}
что совпадает с выражениями~\cite{EqNote}: $I_x=-\gamma_x(1-D_1/\rho)-(2\alpha_x/\rho)D_1'$, $I_y=-\gamma_y(1-D_1/\rho)$ в диполе и $I=-\gamma\pm K_2D_1\beta$ на квадруполях, секступолях и промежутках. Отдельно:
\begin{align}
\label{eq:xiA}
\xi^{\mathrm{A}}_{x,y}&=-\frac{1}{4\pi}\oint\gamma_{x,y}\,\mathrm{d}s
=-\frac{1}{4\pi}\oint K_{x,y}\beta_{x,y}\,\mathrm{d}s,\\
\label{eq:xiC}
\xi^{\mathrm{C}}_{x,y}&=+\frac{1}{4\pi}\int_{\text{диполи}}\frac{D_1\gamma_{x,y}}{\rho}\,\mathrm{d}s,\\
\label{eq:xiE}
\xi^{\mathrm{E}}_{x}&=-\frac{1}{2\pi}\int_{\text{диполи}}\frac{D_1'\alpha_x}{\rho}\,\mathrm{d}s,
\qquad\xi^{\mathrm{E}}_y=0,\\
\label{eq:xiS}
\xi^{\mathrm{S}}_{x,y}&=\pm\frac{1}{4\pi}\oint K_2D_1\beta_{x,y}\,\mathrm{d}s.
\end{align}
Во второй форме \eqref{eq:xiA} использовано $K\beta=\gamma+\alpha'$ (следует из уравнения для бетатронной функции) и $\oint\alpha'\,\mathrm{d}s=0$. Эта замена законна только для интеграла по всему кольцу, причём $K_x$ включает слабую фокусировку $1/\rho^2$ диполей.

\paragraph{Член A: кинематический.}
Он есть в любом элементе, в том числе в дрейфе: в канонических переменных фокусировка от $\delta$ не зависит, а скорость поворота $x'=p_x/(1+\delta)$ --- зависит. В форме $-\oint K\beta/4\pi$ он выглядит как <<квадрупольная>> хроматичность $\Delta K=-K\delta$ уравнения \eqref{eq:hill}, но происходит из кинетического члена. По этой причине формулу $\Delta\nu=(1/4\pi)\oint\beta\Delta K$ нельзя применять поэлементно внутри диполя~\cite{EqNote}: там $x'\neq p_x$ уже в линейном по $\delta$ порядке, и члены C, E не сводятся к $\Delta K$.

\paragraph{Член C: кривизна тела.}
Частица на дисперсионной орбите в диполе идёт по большему ($D_1>0$) радиусу, и её путь на единицу $s$ длиннее в $1+D_1\delta/\rho$ раз. Бетатронный набег фазы растёт пропорционально; член одинаков по форме для $x$ и $y$ (в $y$ --- через $\gamma_y$).

\paragraph{Член E: перекрёстный.}
Наклон дисперсионной орбиты $D_1'\delta$ вместе с множителем $1+X/\rho$ даёт член $XP_x$, т.\,е. связь координаты и импульса горизонтальных колебаний. Он только горизонтальный: аналогичный член по $y$ требовал бы вертикальной дисперсии. Его среднее $\propto\alpha_x$, поэтому член чувствителен к тому, где в диполе находится перетяжка $\beta_x$.

\paragraph{Член S: секступоли.}
Секступоль на дисперсионной орбите --- квадруполь с градиентом $K_2D_1\delta$. Знак разный для $x$ и $y$; это основа коррекции хроматичности двумя семействами~\cite{ValetovPhD}.

\paragraph{Сумма тела.}
Сумма \eqref{eq:xiA}--\eqref{eq:xiS} есть хроматичность гамильтониана \eqref{eq:H}, в котором поле диполя обрывается на торцах без продольной компоненты. Такая модель <<только тело>> точна внутри элементов, но поле с разрывом $B_y$ и $B_s=0$ не удовлетворяет $\operatorname{rot}\mathbf{B}=0$ на границе. Недостающий вклад --- предмет следующей главы.

\section{Край магнита}
\label{sec:edge}

Поле реального диполя на торце спадает, и по уравнениям Максвелла при этом появляется продольная компонента $B_s$. В пределе жёсткого края она становится $\delta$-функцией и даёт тонкий вертикальный удар, пропорциональный горизонтальному наклону траектории. На дисперсионной орбите наклон равен $D_1'\delta$, т.\,е. частица с $\delta\neq0$ входит в магнит с перпендикулярным торцом под эффективным углом. Это обычная краевая фокусировка, но хроматическая. В главе выведен её вклад в $\xi_y$, объяснено, почему в $\xi_x$ его нет, и описан переход от конечного краевого поля к жёсткому краю.

\subsection{Продольное поле жёсткого края}

Пусть вход в диполь находится при $s=0$, торец перпендикулярен орбите, и в средней плоскости $B_y=B_0\Theta(s)$. В вакууме $\operatorname{rot}\mathbf{B}=0$, и $s$-компонента условия $\partial B_s/\partial y=\partial B_y/\partial s$ в первом порядке по $y$ даёт
\begin{equation}
\label{eq:Bs}
B_s=y\,B_0\,\delta_D(s)\quad\text{(вход)},
\qquad
B_s=-y\,B_0\,\delta_D(s-L)\quad\text{(выход)},
\end{equation}
где $\delta_D$ --- дельта-функция Дирака, $L$ --- длина магнита. Кривизной орбиты на бесконечно тонком слое можно пренебречь. Гамильтониан \eqref{eq:H} этой компоненты не содержит: он построен для поля, не зависящего от $s$.

\subsection{Краевой удар}

Сила Лоренца $\mathrm{d}\mathbf{P}/\mathrm{d}t=q\,\mathbf{v}\times\mathbf{B}$. Вертикальная компонента $(\mathbf{v}\times\mathbf{B})_y=v_sB_x-v_xB_s$; в слое $B_x$ не имеет особенности, остаётся член с $B_s$. Интегрирование по слою ($\mathrm{d}t=\mathrm{d}s/v_s$, $v_x/v_s=x'$) и деление на $p_s$ дают
\begin{equation}
\label{eq:kick}
\Delta p_y=-\frac{x'\,y}{\rho}\quad\text{(вход)},
\qquad
\Delta p_y=+\frac{x'\,y}{\rho}\quad\text{(выход)},
\end{equation}
с $qB_0/p_s=1/\rho$. Горизонтальная компонента $(\mathbf{v}\times\mathbf{B})_x=v_yB_s-v_sB_y$ даёт от $B_s$ удар $\propto y'y/\rho$: он квадратичен по вертикальным переменным и от $X$ не зависит.

Это та же формула, что для торца, повёрнутого на угол $e$: вертикальная краевая фокусировка $\Delta p_y=-(\tan e/\rho)\,y$. Для перпендикулярного торца роль $\tan e$ играет наклон траектории $x'$: на входе эффективный угол $e+x'$, на выходе $e-x'$ (в принятой здесь конвенции знаков, в которой $e>0$ фокусирует по $y$). На референсной орбите $x'=0$, и линейной линзы нет. Частица с $\delta\neq0$ идёт по дисперсионной орбите с $x'=D_1'\delta$ и видит тонкую вертикальную линзу.

\subsection{Вклад в хроматичность}

Подстановка $x'=D_1'\delta$ в \eqref{eq:kick} даёт на входе $\Delta p_y=-k\,y$ с $k=D_1'\delta/\rho$ (фокусировка), на выходе $k=-D_1'\delta/\rho$. Тонкая линза сдвигает число колебаний на $\Delta\nu_y=\beta_yk/(4\pi)$, и
\begin{equation}
\label{eq:xiF}
\xi^{\mathrm{F}}_y=\frac{1}{4\pi\rho}\sum_{\text{диполи}}
\Bigl[\beta_yD_1'\big|_{\text{вход}}-\beta_yD_1'\big|_{\text{выход}}\Bigr],
\qquad
\xi^{\mathrm{F}}_x=0.
\end{equation}
В гамильтоновом виде это член $\mathcal{H}_2^{(1)}=\pm(D_1'/2\rho)\,Y^2\,\delta_D(s-s_{\text{торца}})$, который входит в \eqref{eq:xi-ham} так же, как остальные (раздел~\ref{sec:def}). Если оптика симметрична относительно середины диполя, $D_1'$ на входе и выходе имеют разные знаки, и вклады торцов складываются, а не сокращаются. Знак суммы определяется оптикой; при $D_1'>0$ на входе и $D_1'<0$ на выходе вклад положителен, т.\,е. уменьшает модуль отрицательной натуральной $\xi_y$.

\paragraph{Почему нет горизонтального вклада.}
Обычная горизонтальная краевая фокусировка $+(\tan e/\rho)\,x$ возникает потому, что при повёрнутом торце длина пути в поле зависит от $x$. При перпендикулярном торце граница поля --- это поверхность $s=\mathrm{const}$ сопровождающей системы, и путь в поле от $x$ не зависит; наклонная траектория описывается гамильтонианом тела \eqref{eq:H} точно, а её вклад уже содержится в членах C и E. Горизонтальная часть удара от $B_s$ \eqref{eq:kick} квадратична по $y$ и на $\xi_x$ в первом порядке не влияет.

\paragraph{Если торцы повёрнуты.}
При $e\neq0$ к \eqref{eq:xiF} добавляются хроматические части обычной краевой фокусировки: удары $\mp(\tan e/\rho)\,x$, $\pm(\tan e/\rho)\,y$ в канонических импульсах от $\delta$ не зависят, но к ним прибавляются члены от сдвига точки линеаризации (квадратичные коэффициенты края, умноженные на $D_1$) и от зависимости эффективного угла от $D_1'\delta$. Для магнитов этой заметки $e=0$, и подробно этот случай не разбирается.

\subsection{Конечное краевое поле и предел жёсткого края}

Реальный спад описывают функцией Энге~\cite{Enge1967,ValetovPhD}
\begin{equation}
\label{eq:enge}
B_y(s)=B_0F(s),
\qquad
F(s)=\frac{1}{1+\exp\bigl(a_1+a_2(s/g)+\dots+a_n(s/g)^n\bigr)},
\end{equation}
где $g$ --- апертура (зазор), задающая масштаб спада. Положение эквивалентного жёсткого края --- эффективная граница поля (EFB) --- выбирается из равенства площадей:
\begin{equation}
\label{eq:efb}
\int_{-\infty}^{\infty}\bigl[F(s)-\Theta(s_{\mathrm{EFB}}-s)\bigr]\,\mathrm{d}s=0
\end{equation}
(для выхода; для входа зеркально). Тогда интеграл поля по краю совпадает с жёстким краем в $s_{\mathrm{EFB}}$, и угол поворота сохраняется. В~\cite{ValetovPhD} то же условие записано через точки $z_{\mathrm{int}}$, $z_{\mathrm{ext}}$ внутри и вне элемента.

В том же порядке по $y$ продольное поле $B_s=yB_0F'(s)$, и вертикальный удар края
\begin{equation}
\label{eq:kick-fringe}
\Delta p_y=-\frac{1}{\rho}\int x'(s)\,y(s)\,F'(s)\,\mathrm{d}s .
\end{equation}
Если на длине спада $\sim g$ наклон $x'$ и координата $y$ почти постоянны, $\int F'\,\mathrm{d}s=\pm1$ и \eqref{eq:kick-fringe} переходит в \eqref{eq:kick}, а вклад в хроматичность --- в \eqref{eq:xiF} с $\beta_y$, $D_1'$ на EFB. Поправки к пределу жёсткого края имеют порядок $g$:
\begin{itemize}
\item наклон $x'$ и $y$ меняются внутри слоя (кривизна $F/\rho$ действует уже в краевом поле);
\item следующий порядок по $y$ даёт известную поправку к вертикальной краевой фокусировке $\propto g\,K_{\mathrm{fr}}/\rho^2$ (при $e=0$ --- дефокусирующую), где $K_{\mathrm{fr}}=\int F(1-F)\,\mathrm{d}s/g$ (параметр \texttt{FINT} в MAD-X); она сдвигает $\nu_y$ уже при $\delta=0$, а через $1/(1+\delta)$ даёт и малую хроматическую поправку.
\end{itemize}
При $g\to0$ все они исчезают линейно по $g$, и модель с конечным краем стремится к жёсткому краю, но не к модели <<только тело>>: удар \eqref{eq:kick} от формы спада не зависит, пока спад короткий. В~\cite{ValetovChrom2026} это видно по иерархии моделей кольца Muon g-2: жёсткий край на EFB и реалистичный спад Энге дают практически одинаковую хроматичность, а неучёт смещения EFB заметно её меняет.

Отметим, что связь бетатронного удлинения орбиты с хроматичностью $\Lambda=-(2\pi/C_s)(J_x\xi_x+J_y\xi_y)$~\cite{EqNote} от формы края не зависит: обе её части берутся из одного гамильтониана, и член \eqref{eq:xiF} входит в $\Lambda$ через полное $\xi_y$.

\section{Как это считают программы}
\label{sec:codes}

Все программы вычисляют одну величину --- производную \eqref{eq:xi-trace} или эквивалентное ей среднее \eqref{eq:xi-ham}, --- но в разных моделях элементов и разными путями: конечной разностью по $\delta$, продифференцированной картой с энергией-параметром или нормальной формой. Различия результатов определяются не методом, а тем, какие члены главы~\ref{sec:terms} и~\ref{sec:edge} содержит модель диполя. В этой главе разобраны BELL, COSY INFINITY, MAD-X/PTC и OptiM.

\subsection{Карты Тейлора и Ли (BELL)}

Одно-оборотное отображение строится как композиция карт элементов; карта диполя получается из гамильтониана тела \eqref{eq:H} (как ряд Тейлора или как Ли-фактор $\exp(:f:)$), без краевых отображений. Переменная --- канонический $\delta=\delta_p$. Хроматичность берётся прямо по определению: при фиксированном $\delta$ находится неподвижная точка \eqref{eq:fixed}, вычисляются дробные числа колебаний якобиана \eqref{eq:nu-trace}, и $\nu(\delta)$ дифференцируется центральной разностью. Высшие $\nu^{(j)}$ получаются так же из карты высокого порядка. Модель содержит члены A, C, E, S и не содержит F. Нормальная форма 6D (Драгт--Финн) рассчитана на кольцо с ВЧ и для $\xi$ в 4D не используется.

\subsection{Дифференциальная алгебра (COSY INFINITY)}

COSY строит карты в дифференциальной алгебре (DA): каждая координата --- усечённый многочлен от начальных координат и, возможно, от параметров~\cite{Berz1999,COSYManual}. Команда \texttt{OV}~$n$~$n_v$~$n_p$ задаёт порядок $n$, число фазовых пар $n_v$ и число параметров $n_p$. Переменные COSY --- $(x,a,y,b,l,\delta_K)$ с $a=p_x/p_0$, $b=p_y/p_0$~\cite{ValetovPhD}, поэтому сырые хроматичности относятся к $\delta_K$ и переводятся по \eqref{eq:xi-conv}.

\paragraph{Энергия как параметр: \texttt{OV 3 2 1}.}
Кинетическая энергия референсной частицы задаётся как $K_s\cdot\mathrm{PARA}(1)$, т.\,е. $K=K_s(1+\varepsilon)$ с DA-переменной $\varepsilon\equiv\delta_K$. Карта четырёхмерна, но её коэффициенты --- многочлены от $\varepsilon$. Нормальная форма (\texttt{TP}) поворачивает две поперечные пары, и числа колебаний получаются как функции параметра: $\nu(\varepsilon)=\nu_0+(\mathrm{d}\nu/\mathrm{d}\delta_K)\,\varepsilon+\dots$ Это в точности определение \eqref{eq:nu-series}: неподвижная точка при данном $\varepsilon$ находится в DA автоматически, и второе слагаемое \eqref{eq:dM} учтено~\cite{ValetovPhD}. Условие корректности --- поля магнитов должны оставаться фиксированными при изменении $\varepsilon$; если бы градиенты масштабировались с жёсткостью, хроматичность была бы тождественно нулевой.

\paragraph{Энергия как координата: \texttt{OV 3 3 0}.}
Карта шестимерна, $\delta_K$ --- шестая координата, и в ней есть коэффициенты $(x|\delta_K)$, $(x|x\delta_K)$, $(a|a\delta_K)$ и т.\,д. Хроматичность извлекается методом неподвижной точки~\cite{ValetovPhD}: в блоке $(x,a,\delta_K)$ решается $\lambda=(M-I_H)^{-1}e_3$, $I_H=\operatorname{diag}(1,1,0)$, что даёт $(D,D')$; карта пересчитывается около $\mathring z=(D,D',\dots,0,\delta)\delta$, и \eqref{eq:xi-trace} записывается явно через коэффициенты первого и второго порядков (приложение~\ref{app:map}). Результат тот же, что при \texttt{OV 3 2 1}: обе постановки дифференцируют одно и то же семейство отображений.

\paragraph{Почему \texttt{TP} не работает на 6D-карте без ВЧ.}
Нормальная форма приводит линейную часть к трём поворотам. Без ВЧ продольный блок
\begin{equation}
\label{eq:jordan}
M_{l\delta}=\begin{pmatrix}1&(l|\delta_K)\\0&1\end{pmatrix}
\end{equation}
имеет двукратное собственное значение $1$ и не диагонализуется --- это жорданова клетка. Поворота нет, продольного действия нет, и нормальная форма 6D не существует. С ВЧ блок становится поворотом на $2\pi\nu_s$, и \texttt{TP} работает, но $\delta$ уже не параметр, а фазовая переменная: линейная хроматичность в нормальной форме не появляется (глава~\ref{sec:rf}).

\paragraph{Краевые поля: \texttt{FR}.}
Режим \texttt{FR 0} --- жёсткий край без краевого поля: диполь описывается только телом, член F отсутствует, и результат совпадает с BELL. Режимы \texttt{FR 1}--\texttt{FR 3} включают краевое поле со спадом Энге, масштабированным апертурой; \texttt{FR 3} интегрирует уравнения движения в поле, удовлетворяющем уравнениям Максвелла, на всём протяжении края. Эти режимы содержат $B_s$ и, следовательно, член F; при апертуре $\to0$ они стремятся к пределу жёсткого края \eqref{eq:xiF}, а при конечной апертуре отличаются от него на поправки $O(g)$ раздела~\ref{sec:edge}. Для \texttt{FR 1}, \texttt{FR 2} нужна таблица \texttt{SYSCA.DAT}~\cite{COSYManual}.

\subsection{MAD-X и PTC}

\paragraph{MAD-X \texttt{TWISS}.}
Толстый \texttt{SBEND} описывается матрицами TRANSPORT первого и второго порядков~\cite{Brown1968,MADX}, куда входят и краевые матрицы диполя. Хроматичность \texttt{DQ} вычисляется по хроматическим функциям, т.\,е. по формулам вида \eqref{eq:xi-ham} (с флагом \texttt{CHROM} --- конечной разностью по $\delta$), и относится к $p_t$: $\xi=\beta_s\,\mathtt{DQ}$ \eqref{eq:xi-conv}. Член F присутствует даже при $E_1=E_2=0$, $\mathtt{FINT}=0$: краевая матрица второго порядка содержит кинематический член $\propto x'y$. Ненулевые \texttt{FINT}, \texttt{HGAP} добавляют поправку конечного края раздела~\ref{sec:edge} в приближении тонкого слоя.

\paragraph{PTC.}
В точной модели (\texttt{exact}) тело диполя интегрируется по полному корню \eqref{eq:H}, а жёсткий край всегда содержит кинематический краевой член, даже при выключенном \texttt{fringe}; \texttt{fringe=true} добавляет краевое поле. При \texttt{time=false} переменная --- $\delta_p$. Хроматичность и её высшие порядки получаются из нормальной формы с $\delta$ как параметром (\texttt{ptc\_normal}, \texttt{ptc\_twiss})~\cite{PTC}. Приближённая (expanded) модель теряет часть членов с $1/\rho$, а тонкие диполи (\texttt{MAKETHIN}) --- члены C, E и F; для хроматичности их использовать нельзя.

\subsection{OptiM}

OptiM~\cite{OptiM} вычисляет натуральную хроматичность по $\delta_p$ из линейной оптики секторного магнита с жёстким краем, включая краевой член F; результаты совпадают с MAD-X и PTC. Высших хроматичностей и краевого поля конечной длины в используемой модели нет.

\subsection{Python: точный гамильтониан}

Независимая проверка --- прямое интегрирование уравнений Гамильтона для \eqref{eq:H} в теле элементов, поиск замкнутой орбиты Ньютоном при фиксированном $\delta$ и вычисление \eqref{eq:xi-trace} по якобиану. Без краевых ударов результат воспроизводит модель <<только тело>> (BELL, COSY \texttt{FR 0}); с ударами \eqref{eq:kick} на торцах --- MAD-X, PTC и OptiM. Удары обратного знака не воспроизводят ни одну программу, что подтверждает знак \eqref{eq:kick}.

\section{Хроматичность при синхротронном движении}
\label{sec:rf}

С ВЧ $\delta$ перестаёт быть параметром: она колеблется с синхротронной частотой, малой по сравнению с бетатронной. Бетатронная фаза набирается с мгновенной частотой $\nu(\delta(t))$, и наблюдаемое число колебаний есть её среднее по синхротронному периоду. Линейная хроматичность при усреднении выпадает, остаётся член второго порядка. Поэтому нормальная форма 6D (\texttt{TS} в COSY) даёт не $\xi$, а усреднённую вторую хроматичность. Линейную $\xi$ можно восстановить, сдвинув равновесную энергию --- неподвижную точку 6D-отображения --- смещением частоты ВЧ. Глава связывает это с равновесным сдвигом~\cite{EqNote}.

\subsection{Медленная переменная и усреднение}

Продольное движение (см.~\cite{EqNote}, глава~1) при $\nu_s\ll\nu_{x,y}$ даёт
\begin{equation}
\label{eq:delta-t}
\delta(t)=\delta_\eq+\delta_m\cos\psi_s+O(\delta_m^2),
\qquad
\langle\delta\rangle=\delta_\eq,
\qquad
\langle\delta^2\rangle=\delta_\eq^2+\frac{\delta_m^2}{2}.
\end{equation}
За один бетатронный период $\delta$ почти не меняется, и бетатронное движение адиабатически следует за замкнутой орбитой $x_{\mathrm{co}}(s;\delta(t))$. Бетатронная фаза за $N$ оборотов равна $2\pi\sum_k\nu(\delta_k)$, поэтому средняя частота --- это среднее мгновенной (метод усреднения~\cite{Bogoliubov1961}):
\begin{equation}
\label{eq:nu-avg}
\bar\nu=\bigl\langle\nu(\delta)\bigr\rangle
=\nu_0+\xi\,\delta_\eq+\frac{\nu^{(2)}}{2}\Bigl(\delta_\eq^2+\frac{\delta_m^2}{2}\Bigr)+\dots
\end{equation}
Модуляция частоты с синхротронной частотой даёт боковые полосы $\nu\pm k\nu_s$, но среднюю частоту не сдвигает: фаза есть интеграл частоты, и колеблющаяся часть $\xi\delta_m\cos\psi_s$ в среднее не входит. Для частицы около синхронной ($\delta_\eq\approx0$) линейная хроматичность из $\bar\nu$ выпадает; остаётся амплитудная зависимость от синхротронного действия через $\nu^{(2)}$.

\subsection{Что даёт нормальная форма 6D}

В 6D нормальной форме гамильтониан в переменных действие--угол зависит только от действий,
\begin{equation}
\label{eq:nf6}
\tilde{\mathcal{H}}=\mu_xJ_x+\mu_yJ_y+\mu_sJ_s
+\tfrac12\sum_{i,k}c_{ik}J_iJ_k+\dots,
\qquad
\nu_x(J)=\frac{1}{2\pi}\frac{\partial\tilde{\mathcal{H}}}{\partial J_x}.
\end{equation}
Синхротронное действие связано с амплитудой: $\langle\delta^2\rangle_{\psi_s}=\gamma_LJ_s$, где $\gamma_L$ --- продольная функция Твисса для $\delta$. Сравнение с \eqref{eq:nu-avg} при $\delta_\eq=0$ даёт
\begin{equation}
\label{eq:dnu-dJs}
\frac{\partial\nu_x}{\partial J_s}=\frac{\nu_x^{(2)}}{2}\,\gamma_L,
\end{equation}
т.\,е. перекрёстный амплитудный сдвиг <<бетатрон--синхротрон>>, который выдаёт \texttt{TS}, пропорционален второй хроматичности. Линейная $\xi$ в \eqref{eq:nf6} не появляется: в 6D нет параметра, по которому её можно было бы взять, а член $\xi\,\delta\,\mathcal{H}$ с колеблющимся $\delta$ нормальная форма убирает как нерезонансный. При $V\to0$ величина \eqref{eq:dnu-dJs} в нормировке COSY стремится к нулю вместе с $\nu_s$, а не к $\xi$. Если в 6D-постановке задать энергию ещё и параметром (\texttt{OV 3 3 1}), параметр в отображение не попадает: энергия уже есть координата, и производная по параметру равна нулю.

\subsection{Восстановление $\xi$ сдвигом неподвижной точки}

Сместим частоту ВЧ: $\omega_\rf\to\omega_\rf(1+\epsilon)$. Синхронной становится частица, у которой период обращения короче в $1+\epsilon$ раз, т.\,е. $\tau(\delta^*)=-\epsilon$, и в линейном порядке
\begin{equation}
\label{eq:dstar}
\delta^*=-\frac{\epsilon}{\eta_0}.
\end{equation}
Неподвижная точка 6D-отображения сдвигается в $(x,p_x)=(D_1,D_1')\delta^*$, $\delta=\delta^*$. Её якобиан $6\times6$ имеет три пары собственных значений $e^{\pm i\mu}$; поперечные частоты $\nu_{x,y}(\delta^*)$ считаются на смещённой орбите, и
\begin{equation}
\label{eq:xi-fp}
\xi=\frac{\mathrm{d}\nu_{x,y}}{\mathrm{d}\delta^*}\Big|_{\delta^*=0}.
\end{equation}
Это то же определение \eqref{eq:xi-trace}, только неподвижная точка находится в 6D, а $\delta$ задаётся частотой ВЧ. Синхротронная связь через дисперсию меняет собственные поперечные частоты якобиана лишь в высших порядках по малой связи, поэтому \eqref{eq:xi-fp} от напряжения практически не зависит и при $V\to0$ переходит в 4D-результат \texttt{OV 3 2 1}. Тот же приём --- измерение $\xi$ по сдвигу частот при изменении частоты ВЧ --- используется на действующих кольцах.

\subsection{Связь с равновесным сдвигом}

В~\cite{EqNote} показано, что усреднённое по бетатронной фазе удлинение пути равно
\begin{equation}
\label{eq:Lambda}
\Lambda=-\frac{2\pi}{C_s}\bigl(J_x\xi_x+J_y\xi_y\bigr),
\qquad
\delta_\eq=-\frac{1}{\eta_0}\Bigl(\frac{\eta_1\delta_m^2}{2}+\Lambda\Bigr).
\end{equation}
Хроматичность входит сюда как полная, со всеми элементами и краями: вывод~\cite{EqNote} опирается только на гамильтонову форму и статичность поля. Поэтому для $\delta_\eq$ нужно брать $\xi_y$ с членом F \eqref{eq:xiF}; модель <<только тело>> даёт в вертикальной плоскости слишком отрицательную $\xi_y$ и, следовательно, завышенное $\Lambda$. С другой стороны, \eqref{eq:nu-avg} показывает обратную связь: сдвиг $\delta_\eq$ \eqref{eq:Lambda} даёт сдвиг средней частоты $\xi\,\delta_\eq$, т.\,е. член второго порядка по амплитудам, того же порядка, что $\nu^{(2)}\delta_m^2/4$.

\section{Итог: одна физика, разные модели}
\label{sec:summary}

Хроматичность кольца с секторными диполями и перпендикулярными торцами складывается из пяти членов: A (кинематический, $-\oint\gamma/4\pi$), C (кривизна тела, $+\int D_1\gamma/4\pi\rho$), E (перекрёстный, $-\int D_1'\alpha_x/2\pi\rho$, только $x$), S (секступоли, $\pm\oint K_2D_1\beta/4\pi$) и F (жёсткий край, $\sum[\beta_yD_1']^{\text{вход}}_{\text{выход}}/4\pi\rho$, только $y$). Члены A, C, E, S следуют из гамильтониана тела; F --- из продольного поля $B_s$, которого требует $\operatorname{rot}\mathbf{B}=0$ на границе поля. Разные программы дают разные $\xi_y$ не из-за метода вычисления (все методы главы~\ref{sec:codes} эквивалентны определению главы~\ref{sec:def}), а из-за того, есть ли член F в их модели диполя. Горизонтальная хроматичность от модели края в первом порядке не зависит.

\begin{center}
\small
\renewcommand{\arraystretch}{1.2}
\begin{tabular}{@{}lcccccc@{}}
\hline
Модель & A & C & E & S & F & конечн. край\\
\hline
BELL (Тейлор / Ли) & + & + & + & + & --- & ---\\
COSY \texttt{FR 0} & + & + & + & + & --- & ---\\
COSY \texttt{FR 1}--\texttt{FR 3} & + & + & + & + & + & Энге\\
MAD-X \texttt{TWISS}, \texttt{SBEND} & + & + & + & + & + & \texttt{FINT}, \texttt{HGAP}\\
PTC, \texttt{exact} & + & + & + & + & + & \texttt{fringe=true}\\
OptiM & + & + & + & + & + & ---\\
Python, тело & + & + & + & + & --- & ---\\
Python, тело + \eqref{eq:kick} & + & + & + & + & + & ---\\
\hline
\end{tabular}
\end{center}

Замечания к таблице.
\begin{itemize}
\item Строки <<Python>> --- независимая проверка интегрированием точного гамильтониана \eqref{eq:H} (глава~\ref{sec:codes}).
\item <<Конечный край>> означает поправки $O(g)$ раздела~\ref{sec:edge}; при $g\to0$ строки COSY \texttt{FR 1}--\texttt{FR 3} переходят в строку жёсткого края (MAD-X, PTC, OptiM), а не в строку \texttt{FR 0}.
\item Физически правильная модель диполя с перпендикулярными торцами и коротким спадом поля --- с членом F. Модель <<только тело>> соответствует полю, не удовлетворяющему уравнениям Максвелла на границе.
\item Для COSY с энергией-параметром (\texttt{OV 3 2 1}) и с энергией-координатой (\texttt{OV 3 3 0}, метод неподвижной точки) результат один и тот же; переменная COSY --- $\delta_K$, перевод \eqref{eq:xi-conv}. Для MAD-X \texttt{DQ} относится к $p_t$.
\item С ВЧ линейная $\xi$ получается только сдвигом неподвижной точки \eqref{eq:xi-fp}; \texttt{TS} даёт усреднённую вторую хроматичность \eqref{eq:dnu-dJs}.
\item В $\Lambda$ и равновесный сдвиг \eqref{eq:Lambda} входит полная хроматичность, с членом F.
\item Приближённые модели --- PTC без \texttt{exact}, тонкие диполи после \texttt{MAKETHIN} --- теряют члены с $1/\rho$ (C, E, F) и в таблицу не включены.
\end{itemize}


\clearpage
\appendix
\section{Хроматичность через коэффициенты отображения}
\label{app:map}

Запишем формулу \eqref{eq:xi-trace} явно через коэффициенты одно-оборотного отображения, как это сделано в~\cite{ValetovPhD,ValetovChrom2026}, но в переменных этой заметки $(x,a,y,b,\delta)$, $a=p_x$, $b=p_y$, $\delta=\delta_p$. Коэффициент $(x|x^ia^j\delta^k)$ --- множитель при $x^ia^j\delta^k$ в разложении конечного $x$ по начальным переменным, и так же для остальных. Кольцо плоское: коэффициенты, нечётные по $(y,b)$, в уравнениях для $x$, $a$ равны нулю.

\paragraph{Неподвижная точка.}
В первом порядке по $\delta$ неподвижная точка $(x,a)=(D,D')\delta$ удовлетворяет
\begin{equation}
\label{eq:app-fp}
\begin{pmatrix}1-(x|x)&-(x|a)\\-(a|x)&1-(a|a)\end{pmatrix}
\begin{pmatrix}D\\D'\end{pmatrix}
=\begin{pmatrix}(x|\delta)\\(a|\delta)\end{pmatrix},
\end{equation}
откуда, с учётом $(x|x)(a|a)-(x|a)(a|x)=1$,
\begin{equation}
\label{eq:app-D}
D=\frac{\bigl(1-(a|a)\bigr)(x|\delta)+(x|a)(a|\delta)}{2-2\cos\mu_x},
\qquad
D'=\frac{(a|x)(x|\delta)+\bigl(1-(x|x)\bigr)(a|\delta)}{2-2\cos\mu_x}.
\end{equation}
Это значения $D_1$, $D_1'$ в точке наблюдения.

\paragraph{Матрица на смещённой орбите.}
Диагональные элементы якобиана в точке $(D,D')\delta$ в первом порядке по $\delta$:
\begin{align}
\tilde M_{11}&=(x|x)+\delta\bigl[2(x|xx)D+(x|xa)D'+(x|x\delta)\bigr],\\
\tilde M_{22}&=(a|a)+\delta\bigl[(a|xa)D+2(a|aa)D'+(a|a\delta)\bigr],\\
\tilde M_{33}&=(y|y)+\delta\bigl[(y|xy)D+(y|ay)D'+(y|y\delta)\bigr],\\
\tilde M_{44}&=(b|b)+\delta\bigl[(b|xb)D+(b|ab)D'+(b|b\delta)\bigr].
\end{align}
Внедиагональные элементы в след не входят.

\paragraph{Хроматичность.}
По \eqref{eq:xi-trace}
\begin{equation}
\label{eq:app-xix}
\xi_x=-\frac{
2(x|xx)D+(x|xa)D'+(x|x\delta)+(a|xa)D+2(a|aa)D'+(a|a\delta)
}{4\pi\sin\mu_x},
\end{equation}
\begin{equation}
\label{eq:app-xiy}
\xi_y=-\frac{
(y|xy)D+(y|ay)D'+(y|y\delta)+(b|xb)D+(b|ab)D'+(b|b\delta)
}{4\pi\sin\mu_y},
\end{equation}
где $\cos\mu=\tfrac12\operatorname{tr}M$, а знак $\sin\mu$ совпадает со знаком $(x|a)$ (соответственно $(y|b)$). Подстановка \eqref{eq:app-D} даёт формулу, содержащую только коэффициенты первого и второго порядков; это та же структура, что в~\cite{ValetovPhD} (там $\delta_K$ и обозначения $M^j_{i_1\dots i_6}$).

\paragraph{Где здесь члены A--F.}
Коэффициенты $(x|x\delta)$, $(a|a\delta)$, $(y|y\delta)$, $(b|b\delta)$ --- прямая зависимость от $\delta$ (член A и та часть C, E, что не требует сдвига орбиты), а произведения квадратичных коэффициентов на $D$, $D'$ --- второе слагаемое \eqref{eq:dM}, т.\,е. секступоли и геометрическая нелинейность тела диполя. Краевой удар \eqref{eq:kick} на входе добавляет к карте торца коэффициент $(b|ay)=-1/\rho$, на выходе $+1/\rho$. На смещённой орбите $a=D_1'\delta$ он даёт элемент $\partial b/\partial y=\mp D_1'\delta/\rho$ карты торца; после композиции с остальным кольцом он переходит в коэффициенты $(y|ay)$, $(b|ab)$, $(y|y\delta)$, $(b|b\delta)$ одно-оборотного отображения и даёт член F. В модели <<только тело>> этого коэффициента в картах торцов нет.


\clearpage
\bibliographystyle{unsrt}
\bibliography{references}

\end{document}
